\section{Данные}

\subsection{Датасет SciERC}

SciERC~\cite{luan2018scierc} --- корпус научных абстрактов из 12 конференций по ИИ (ACL, EMNLP, NeurIPS и др.). Датасет содержит 500 документов, размеченных на уровне именованных сущностей, отношений и кореференций. В данной работе используется только задача NER.

Стандартное разбиение: 350 документов в обучающей выборке, 50 в валидационной, 100 в тестовой.

\subsection{Типы сущностей}

Датасет определяет 6 типов именованных сущностей (таблица~\ref{tab:entity_types}).

\begin{table}[h]
\centering
\caption{Типы сущностей в SciERC}
\label{tab:entity_types}
\begin{tabular}{lll}
\hline
\textbf{Тип} & \textbf{Описание} & \textbf{Пример} \\
\hline
Task               & Задача или проблема              & \textit{machine translation} \\
Method             & Метод, модель, система           & \textit{BERT, attention mechanism} \\
Metric             & Метрика оценки                   & \textit{F1 score, BLEU} \\
Material           & Датасет, ресурс, материал        & \textit{Penn Treebank, ImageNet} \\
OtherScientificTerm & Прочий научный термин           & \textit{feature, embedding} \\
Generic            & Общий термин, сложно категоризировать & \textit{approach, model} \\
\hline
\end{tabular}
\end{table}

\subsection{BIO-разметка}

Сущности кодируются в формате BIO (Beginning--Inside--Outside):
\begin{itemize}
    \item \texttt{B-X} --- первый токен сущности типа X,
    \item \texttt{I-X} --- продолжение сущности типа X,
    \item \texttt{O} --- токен вне именованной сущности.
\end{itemize}

Итого 13 меток: \texttt{O} плюс \texttt{B-}/\texttt{I-} для каждого из 6 типов. Классы существенно несбалансированы: метка \texttt{O} составляет большинство токенов, а такие типы, как Metric и Task, встречаются значительно реже остальных.

\subsection{Статистика}

Дисбаланс классов --- ключевая особенность датасета, влияющая на выбор функции потерь. В разделе~\ref{sec:experiments} описано, как это учитывается при обучении.
